\subsection{换环}

设 A 与 B 是环，\(f : A \to B\) 是一个环同态。设 C 是 A-模的一个加性集，D 是 B-模的一个加性集。

首先，假设对每个 D 型 B-模 M，将标量环限制到 A 得到的 A-模 \(f_{*}(M)\) 是 C 型的。则从 D 到 \(\mathrm{K}(\mathcal{C})\) 的、将 M 映到 \([f_{*}(M)]_{\mathcal{C}}\) 的映射是模的一个加性函数；由它我们导出一个群同态

\[
f _ {*}: \mathrm{K} (\mathcal {D}) \longrightarrow \mathrm{K} (\mathcal {C}).
\]

我们同样定义一个群同态 \(f_{*}:\mathrm{K}^{\prime}(\mathcal{D})\to \mathrm{K}^{\prime}(\mathcal{C})\)。

我们现在假设对每个 C 型 A-模 E，经由 f 作标量扩张从 E 得到的 B-模 \(f^{*}(\mathrm{E})\)（II, §5, No.~1, 第 277 页）是 D 型的。从 C 到 \(\mathrm{K}'(\mathcal{D})\) 的、将 C 的元素 E 映到 \([f^{*}(\mathrm{E})]_{\mathcal{D}}'\) 的映射是模的一个弱加性函数；因此它诱导一个群同态 \(f^{*}: \mathrm{K}'(\mathcal{C}) \to \mathrm{K}'(\mathcal{D})\)。

再假设对每个正合列

\[
0 \longrightarrow \mathrm{E} ^ {\prime} \longrightarrow \mathrm{E} \longrightarrow \mathrm{E} ^ {\prime \prime} \longrightarrow 0
\]

（由 C 型 A-模组成），D 型 B-模的序列

\[
0 \longrightarrow \mathrm{B} \otimes_ {\mathrm{A}} \mathrm{E} ^ {\prime} \longrightarrow \mathrm{B} \otimes_ {\mathrm{A}} \mathrm{E} \longrightarrow \mathrm{B} \otimes_ {\mathrm{A}} \mathrm{E} ^ {\prime \prime} \longrightarrow 0
\]

都是正合的。这在下述情形特别成立：

\begin{enumerate}
\def\labelenumi{\alph{enumi})}
\item
  同态 \(f\) 使 B 成为一个投射 A-模（II, §3, No.~6, 第 251 页，命题 5 与 II, §3, No.~7, 第 257 页，推论 6）*，或更一般地，一个平坦（X, §1, n° 3, 第 8 页, définition 1）A-模*。
\item
  *集 C 是投射 A-模的类所成之集，或更一般地，平坦（X, §4, n° 5, 第 72 页, corollaire 2）A-模的类所成之集*。
\end{enumerate}

这时从 \(\mathcal{C}\) 到 \(\mathrm{K}(\mathcal{D})\) 的、将 \(\operatorname{E}\) 映到 \([f^{*}(\operatorname{E})]_{\mathcal{D}}\) 的映射是加性的。因此它诱导一个群同态 \(f^{*}:\mathrm{K}(\mathcal{C})\to \mathrm{K}(\mathcal{D})\)。
% semantic-fragments: legacy-input-seam
\relax{}
